\documentclass[11pt,a4paper]{article}
\usepackage[T1]{fontenc}
\usepackage{lmodern,amsmath,amssymb,geometry}
\geometry{margin=30mm}
\title{EPPA--III Lemma 3.12, Case 2: a direct degree argument}
\author{Source-local validation supplement}
\date{9 October 2026}
\begin{document}
\maketitle

\noindent
\textbf{Scope.} This replaces only the sentence in
Case 2 that refers to Claims 3.5 and 3.6
for the double-cover conclusion. The main manuscript
is untouched. We assume that the preceding paragraph
has correctly established \(n\) invariant fibre blocks
of size two with exactly one point of \(G\) in each.

\paragraph{Proof.}
Write the fibres as \(\{x_0,x_1\}\)
for \(x\in G=\dot\cup_i K_{k_i}\).
The internal relation \(\varepsilon\in\{0,1\}\)
between \(x_0\) and \(x_1\) is independent of
\(x\), because \(H\) is vertex-transitive and
the system of pairs is invariant.
Every isomorphism between two ordered edges of
\(G\) extends to an automorphism of \(H\);
the same holds for ordered nonedges.
Such an extension carries the mate of each
domain vertex to the mate of its image.
Therefore, for distinct \(x,y\), all four
adjacencies between the two fibres depend only
on whether \(x\sim_G y\). Reversing the two
endpoints shows that the two off-diagonal
cross-layer entries agree.

Let \(c_E,c_N\in\{0,1\}\) denote these
cross-layer entries above \(G\)-edges and
\(G\)-nonedges, respectively.
Let \(e_T,n_T\in\{0,1\}\) denote the
top-layer entries above \(G\)-edges and
\(G\)-nonedges. For a bottom vertex \(x_0\)
above a clique of size \(k\), its degree is
\[
 d_0(k)=k-1+\varepsilon+c_E(k-1)+c_N(n-k).
\]
All vertices of \(H\) have equal degree.
Since the clique sizes are unequal, comparing
two different values of \(k\) gives
\(1+c_E-c_N=0\). Thus \(c_E=0\) and
\(c_N=1\). In particular,
\[
                       d_0(k)=n-1+\varepsilon.
\]
The top-layer vertex \(x_1\) has degree
\[
 d_1(k)=\varepsilon+e_T(k-1)+n_T(n-k)+(n-k).
\]
As \(d_1(k)=d_0(k)\) for two distinct
clique sizes, subtraction gives
\(e_T-n_T-1=0\).
Hence \(e_T=1\) and \(n_T=0\).
These four entries are exactly those of
the Taylor double of \(G\), and
\(\varepsilon=1\) adds precisely the
matching between fibre mates.
This proves the desired double-cover
identification.

\paragraph{Finishing Case 2.}
With this explicit identification,
the previously established triangle formula
\[
 T_i-T_j=(k_i-k_j)(n-k_i-k_j)
\]
excludes three or more cliques of unequal sizes,
by vertex-transitivity. Hence \(s=2\).

\paragraph{Verification.}
The adjacency-matrix degree argument is
independently checked for all 7,520
cases obtained by choosing one of 32 possible
matrix patterns and an unequal positive
composition of \(n\), for \(3\le n\le8\).
Only the two expected Taylor-double patterns
are regular. See
\texttt{checks/section3\_pair\_regular.py}.
This is a complete paper-level proof of
the reduction \emph{conditional on the previously
established invariant pair system};
it is not yet a full Lean formalization.
\end{document}
